\section*{الفصل \ref{ch:b1:periodicity} --- الدورية}

\begin{solution}{exo:b1:periodicity:1}
\ce{Ca} $[\ce{Ar}]\,4\text{s}^2$: \omterm{def:b1:periodicity:table}{الدورة} 4، \omterm{def:b1:periodicity:table}{الزمرة} 2، \omterm{def:b1:periodicity:table}{الفئة} s. \ce{Fe}
$[\ce{Ar}]\,3\text{d}^6 4\text{s}^2$: \omterm{def:b1:periodicity:table}{الدورة} 4، \omterm{def:b1:periodicity:table}{الزمرة} 8، \omterm{def:b1:periodicity:table}{الفئة} d. \ce{Se}
$[\ce{Ar}]\,3\text{d}^{10} 4\text{s}^2 4\text{p}^4$: \omterm{def:b1:periodicity:table}{الدورة} 4، \omterm{def:b1:periodicity:table}{الزمرة} 16، \omterm{def:b1:periodicity:table}{الفئة}
p. \ce{I} $[\ce{Kr}]\,4\text{d}^{10} 5\text{s}^2 5\text{p}^5$: \omterm{def:b1:periodicity:table}{الدورة} 5،
\omterm{def:b1:periodicity:table}{الزمرة} 17، \omterm{def:b1:periodicity:table}{الفئة} p.
\end{solution}

\begin{solution}{exo:b1:periodicity:2}
\ce{Na} أكبر من \ce{Cl}: \omterm{def:b1:periodicity:table}{الدورة} نفسها، و$Z^{*}$ أصغر. و\ce{Na}
أكبر من \ce{Na+}، الذي فقد طبقته 3s. و\ce{Cl-}
(\qty{181}{pm}) أكبر من \ce{F-} (\qty{133}{pm}): طبقة إضافية.
و\ce{Na+} (\qty{102}{pm}) أكبر من \ce{Mg^{2+}} (\qty{72}{pm}): الإلكترونات
العشرة نفسها، وشحنة نووية 11 مقابل 12.
% ledger: rion:Cl-VI, rion:F-VI, rion:Na+VI, rion:Mg2+VI
\end{solution}

\begin{solution}{exo:b1:periodicity:3}
$\ce{Na} (5.14) < \ce{Al} (5.99) < \ce{Mg} (7.65) < \ce{S} (10.36) <
\ce{P} (10.49) < \ce{Ar} (15.76)$. فالتزايد العام على امتداد \omterm{def:b1:periodicity:table}{الدورة}
ينكسر مرتين: \ce{Al} تحت \ce{Mg} (الإلكترون المنزوع هو أول إلكترون
3p، وهو أعلى طاقةً من 3s) و\ce{S} تحت \ce{P} (الإلكترون
المنزوع هو أول إلكترون 3p مقترن).
\end{solution}

\begin{solution}{exo:b1:periodicity:4}
تقيس \omterm{def:b1:periodicity:polarisability}{قابلية الاستقطاب} مدى سهولة تشوّه السحابة الإلكترونية تحت تأثير
مجال كهربائي. و\ce{I-} أكثر قابلية للاستقطاب من \ce{F-}: فهو أكبر
وإلكتروناته الخارجية أبعد عن النواة. و\ce{Cl-} (18
إلكترونًا تمسكها $Z = 17$) أكثر قابلية للاستقطاب بكثير من \ce{Na+} (10
إلكترونات تمسكها $Z = 11$)، شأن الأنيونات مع الكاتيونات.
\end{solution}

\begin{solution}{exo:b1:periodicity:5}
النسب: $18.83/5.99 = 3.1$، $28.45/18.83 = 1.5$، $119.99/28.45 = 4.2$.
تأتي القفزة بعد الإلكترون الثالث: ثلاثة \omterm{def:b1:quantum-numbers:configuration}{إلكترونات تكافؤ}، \omterm{def:b1:periodicity:table}{الزمرة} 13،
\omterm{def:b1:periodicity:table}{الدورة} 3: إنه الألومنيوم. ويكوّن \ce{Al^{3+}}، $1\text{s}^2 2\text{s}^2
2\text{p}^6$.
% ledger: ie:Al, ie2:Al, ie3:Al, ie4:Al
\end{solution}

\begin{solution}{exo:b1:periodicity:6}
يدخل الإلكترون المضاف إلى \ce{F-} الطبقة 2p الصغيرة، حيث
تدفعه بقوة \omterm{def:b1:quantum-numbers:configuration}{إلكترونات التكافؤ} السبعة الموجودة فيها؛ وفي الطبقة 3p
الأكبر للكلور يكون التنافر أصغر، فتتحرر طاقة
أكبر. وفي النيتروجين ($2\text{p}^3$، \omorbs{u,u,u}) يتعين على الإلكترون الزائد
أن يقترن في \omterm{def:b1:quantum-numbers:orbital}{طبقة فرعية} نصف ممتلئة؛ وفي المغنيسيوم ($3\text{s}^2$)
يتعين عليه أن يبدأ \omterm{def:b1:quantum-numbers:orbital}{الطبقة الفرعية} 3p الأعلى. وفي الحالتين لا يكون
الأنيون أكثر استقرارًا من الذرة والإلكترون الحر.
\end{solution}

\begin{solution}{exo:b1:periodicity:7}
$\chi_M(\ce{Na}) = (5.14 + 0.55)/2 = \qty{2.85}{eV}$؛ $\chi_M(\ce{Cl}) =
(12.97 + 3.61)/2 = \qty{8.29}{eV}$. الفرق كبير: فالزوج الإلكتروني
لرابطة \ce{Na-Cl} يعود كله تقريبًا إلى الكلور، وكلوريد
الصوديوم أيوني، \ce{Na+ Cl-}.
\end{solution}

\begin{solution}{exo:b1:periodicity:8}
\ce{H-F}: $\Delta\chi = 1.78$، على حدود المنطقة الأيونية (وهي
عمليًا رابطة تساهمية شديدة القطبية)، \ce{F^{\delta-}}. \ce{C-H}: 0.35،
شبه غير قطبية. \ce{Na-Cl}: 2.23، أيونية، \ce{Cl^{-}}. \ce{C-Cl}: 0.61،
تساهمية قطبية، \ce{C^{\delta+}-Cl^{\delta-}}. \ce{O-H}: 1.24، تساهمية
قطبية، \ce{O^{\delta-}-H^{\delta+}}.
\end{solution}

\begin{solution}{exo:b1:periodicity:9}
نزولًا في \omterm{def:b1:periodicity:table}{الزمرة} يكون \omterm{def:b1:quantum-numbers:configuration}{إلكترون التكافؤ} في طبقة ذات $n$ أكبر، أبعد
عن النواة وتحجبه \omterm{def:b1:quantum-numbers:configuration}{إلكترونات داخلية} أكثر؛ فيُنزع
بسهولة أكبر. وللروبيديوم، الذي يزيد طبقة، \omterm{def:b1:periodicity:ionisation-energy}{طاقة تأين أولى} أدنى من
\qty{4.34}{eV}.
\end{solution}

\begin{solution}{exo:b1:periodicity:10}
الزنك $[\ce{Ar}]\,3\text{d}^{10} 4\text{s}^2$: يأتي إلكترونه من طبقة
4s ممتلئة؛ أما الغاليوم، $\dots 4\text{s}^2 4\text{p}^1$، فيفقد إلكترونه 4p
الوحيد، الأعلى طاقةً والمحجوب بالطبقتين 4s و3d الممتلئتين. ويفقد الزرنيخ،
$4\text{p}^3$، إلكترونًا من \omterm{def:b1:quantum-numbers:orbital}{طبقة فرعية} نصف ممتلئة؛ أما السيلينيوم،
$4\text{p}^4$، فيفقد الإلكترون المقترن، الذي يرفعه تنافره مع
شريكه: فطاقة تأينه أدنى قليلًا.
\end{solution}

\begin{solution}{exo:b1:periodicity:11}
$\Delta E = (3.98 - 2.20)^2 = 1.78^2 = \qty{3.17}{eV} = 3.17 \times 96.5 =
\qty{306}{kJ/mol}$. ومن ثم $D(\ce{H-F}) = \tfrac12(436 + 158) + 306 = 297 +
306 \approx \qty{603}{kJ/mol}$. والفائض الكبير هو التجاذب الإضافي
بين الشحنتين الجزئيتين $\ce{H^{\delta+}}$ و$\ce{F^{\delta-}}$، الناتج
عن الفرق الكبير في \omterm{def:b1:periodicity:electronegativity}{الكهرسلبية}.
\end{solution}

\begin{solution}{exo:b1:periodicity:12}
يُعرَّف \omterm{def:b1:periodicity:electronegativity}{سلّم بولينغ} من طاقات الروابط A--B: فالعنصر
الذي لا يكوّن أي رابطة ليست له قيمة بولينغ. والهيليوم والنيون والأرغون لا تكوّن
مركبات مستقرة. أما الكريبتون، والزينون خصوصًا، فطبقات تكافئهما أكبر وأكثر
قابلية للاستقطاب وطاقات تأينهما أدنى
(\qty{14.0}{eV} للكريبتون مقابل \qty{15.76}{eV} للأرغون)؛ فيؤكسدهما
الفلور والأكسجين ويكوّنان مركبات مثل \ce{XeF2}، فيمكن قياس
طاقات روابطهما.
% ledger: ie:Kr, ie:Ar
\end{solution}

\begin{solution}{pb:b1:periodicity:1}
\textbf{1.} يُنزع كل إلكترون من أيون أشد موجبيةً وأصغر
من سابقه، فيمسك إلكتروناته بقوة أكبر.
\textbf{2.} $15.04/7.65 = 1.97$؛ $80.14/15.04 = 5.33$؛ $109.27/80.14 = 1.36$.
\textbf{3.} تأتي القفزة بعد إلكترونين: للعنصر X إلكترونا تكافؤ،
فهو في \omterm{def:b1:periodicity:table}{الزمرة} 2؛ وفي \omterm{def:b1:periodicity:table}{الدورة} 3 هو المغنيسيوم.
\textbf{4.} \ce{Mg^{2+}}، $1\text{s}^2 2\text{s}^2 2\text{p}^6$، وهو
توزيع النيون.
\textbf{5.} فلز: \omterm{def:b1:periodicity:table}{الزمرة} 2، في يسار الجدول، بطاقات تأين
أولى منخفضة.
\textbf{6.} يأتي الإلكترون الثالث من قلب 2p، ذي $n$ أصغر،
الأقرب كثيرًا إلى النواة والذي يكاد لا يُحجب، ويُنزع من أيون
ثنائي الشحنة.
% ledger: ie:Mg, ie2:Mg, ie3:Mg, ie4:Mg
\textbf{7.} لكلها عشرة إلكترونات: $1\text{s}^2 2\text{s}^2 2\text{p}^6$.
\textbf{8.} الإلكترونات العشرة نفسها تمسكها شحنات نووية 8 و9 و11
و12 و13: فتنكمش السحابة كلما كبر $Z$.
\textbf{9.} \ce{O^{2-}}: الأكبر، بأدنى شحنة نووية
لإلكتروناته العشرة، وهو أنيون.
\textbf{10.} $140/54 = 2.6$.
\textbf{11.} $198.8/2 = \qty{99.4}{pm}$؛ فالأنيون أكبر بمقدار $181/99.4 = 1.8$
مرة من \omterm{def:b1:periodicity:radius}{نصف القطر التساهمي}.
% ledger: re:Cl2, rion:Cl-VI
\textbf{12.} $\ce{Mg^{2+}} < \ce{Na+} < \ce{Cl-}$: الكاتيونان
متساويا الإلكترونات (شحنة 12 مقابل 11)، وللأيون \ce{Cl-} طبقة ثالثة.
\textbf{13.} $\chi_M$ = 10.41 (\ce{F})، 8.29 (\ce{Cl})، 7.59 (\ce{Br})،
7.53 (\ce{O})، 6.26 (\ce{C})، بالإلكترون فولت.
\textbf{14.} موليكن: $\ce{F} > \ce{Cl} > \ce{Br} > \ce{O} > \ce{C}$؛
بولينغ: $\ce{F} > \ce{O} > \ce{Cl} > \ce{Br} > \ce{C}$. فينتقل الأكسجين
من المرتبة الرابعة إلى الثانية.
\textbf{15.} $3.98/10.41 = 0.382$، $3.16/8.29 = 0.381$، $2.96/7.59 =
0.390$: شبه ثابتة، نحو 0.38؛ فالسلّمان متناسبان في
الهالوجينات.
\textbf{16.} \omterm{def:b1:periodicity:electron-affinity}{الألفة الإلكترونية} للأكسجين صغيرة (\qty{1.44}{eV})
لأن الإلكترون المضاف يجب أن يقترن في الطبقة 2p المتراصّة؛ فقيمة
موليكن للذرة الحرة تبخس الجذب الذي يمارسه الأكسجين في
روابطه، وهو ما يقيسه \omterm{def:b1:periodicity:electronegativity}{سلّم بولينغ}، المبني على الروابط، مباشرة.
\textbf{17.} الكلور في \ce{HCl}؛ والفلور في \ce{ClF}.
\textbf{18.} $\Delta\chi = 0.96$: تساهمية قطبية.
\textbf{19.} $D(\ce{H-H}) = 2 \times 218.0 = \qty{436.0}{kJ/mol}$.
\textbf{20.} $D(\ce{Cl-Cl}) = 2 \times 121.3 = \qty{242.6}{kJ/mol}$.
\textbf{21.} $D(\ce{H-Cl}) = 218.0 + 121.3 - (-92.3) = \qty{431.6}{kJ/mol}$.
\textbf{22.} $\Delta E = 431.6 - \tfrac12(436.0 + 242.6) = 431.6 - 339.3 =
\qty{92.3}{kJ/mol}$. وبالرموز، $D(\ce{H-Cl}) = \Delta_fH(\ce{H}) +
\Delta_fH(\ce{Cl}) - \Delta_fH(\ce{HCl})$ بينما متوسط الرابطتين
المتناظرتين هو $\Delta_fH(\ce{H}) + \Delta_fH(\ce{Cl})$، فيكون $\Delta E =
-\Delta_fH(\ce{HCl})$.
\textbf{23.} $92.3/96.485 = \qty{0.957}{eV}$.
\textbf{24.} $|\chi(\ce{Cl}) - \chi(\ce{H})| = \sqrt{0.957} =
\textbf{0.98}$، مقابل $3.16 - 2.20 = 0.96$ في الجداول: فحساب بولينغ،
مُجرًى بمعطيات حديثة، يعطي الفرق المجدول في حدود
0.02.
\end{solution}
